% Options for packages loaded elsewhere
\PassOptionsToPackage{unicode}{hyperref}
\PassOptionsToPackage{hyphens}{url}
\PassOptionsToPackage{dvipsnames,svgnames,x11names}{xcolor}
\documentclass[
  a4paper,
]{article}
\usepackage{xcolor}
\usepackage[margin=20mm]{geometry}
\usepackage{amsmath,amssymb}
\setcounter{secnumdepth}{-\maxdimen} % remove section numbering
\usepackage{iftex}
\ifPDFTeX
  \usepackage[T1]{fontenc}
  \usepackage[utf8]{inputenc}
  \usepackage{textcomp} % provide euro and other symbols
\else % if luatex or xetex
  \usepackage{unicode-math} % this also loads fontspec
  \defaultfontfeatures{Scale=MatchLowercase}
  \defaultfontfeatures[\rmfamily]{Ligatures=TeX,Scale=1}
\fi
\usepackage{lmodern}
\ifPDFTeX\else
  % xetex/luatex font selection
  \setmainfont[]{DejaVu Serif}
  \setsansfont[]{DejaVu Sans}
  \setmonofont[]{DejaVu Sans Mono}
\fi
% Use upquote if available, for straight quotes in verbatim environments
\IfFileExists{upquote.sty}{\usepackage{upquote}}{}
\IfFileExists{microtype.sty}{% use microtype if available
  \usepackage[]{microtype}
  \UseMicrotypeSet[protrusion]{basicmath} % disable protrusion for tt fonts
}{}
\makeatletter
\@ifundefined{KOMAClassName}{% if non-KOMA class
  \IfFileExists{parskip.sty}{%
    \usepackage{parskip}
  }{% else
    \setlength{\parindent}{0pt}
    \setlength{\parskip}{6pt plus 2pt minus 1pt}}
}{% if KOMA class
  \KOMAoptions{parskip=half}}
\makeatother
\usepackage{longtable,booktabs,array}
\newcounter{none} % for unnumbered tables
\usepackage{calc} % for calculating minipage widths
% Correct order of tables after \paragraph or \subparagraph
\usepackage{etoolbox}
\makeatletter
\patchcmd\longtable{\par}{\if@noskipsec\mbox{}\fi\par}{}{}
\makeatother
% Allow footnotes in longtable head/foot
\IfFileExists{footnotehyper.sty}{\usepackage{footnotehyper}}{\usepackage{footnote}}
\makesavenoteenv{longtable}
\usepackage{graphicx}
\makeatletter
\newsavebox\pandoc@box
\newcommand*\pandocbounded[1]{% scales image to fit in text height/width
  \sbox\pandoc@box{#1}%
  \Gscale@div\@tempa{\textheight}{\dimexpr\ht\pandoc@box+\dp\pandoc@box\relax}%
  \Gscale@div\@tempb{\linewidth}{\wd\pandoc@box}%
  \ifdim\@tempb\p@<\@tempa\p@\let\@tempa\@tempb\fi% select the smaller of both
  \ifdim\@tempa\p@<\p@\scalebox{\@tempa}{\usebox\pandoc@box}%
  \else\usebox{\pandoc@box}%
  \fi%
}
% Set default figure placement to htbp
\def\fps@figure{htbp}
\makeatother
\setlength{\emergencystretch}{3em} % prevent overfull lines
\providecommand{\tightlist}{%
  \setlength{\itemsep}{0pt}\setlength{\parskip}{0pt}}
\usepackage{luatexja}
\usepackage{xurl}
\usepackage{float}
\usepackage{bookmark}
\IfFileExists{xurl.sty}{\usepackage{xurl}
\usepackage{float}}{} % add URL line breaks if available
\urlstyle{same}
\hypersetup{
  pdftitle={Spacetime and Internal-Space Readout Connection Hypothesis in the 3D Triangle-Loop Product (Part III)},
  pdfauthor={Noriaki Kihara (WF System Co., Ltd.) ORCID 0009-0004-6753-4020},
  colorlinks=true,
  linkcolor={Maroon},
  filecolor={Maroon},
  citecolor={Blue},
  urlcolor={Blue},
  pdfcreator={LaTeX via pandoc}}

\title{Spacetime and Internal-Space Readout Connection Hypothesis in the
3D Triangle-Loop Product (Part III)}
\usepackage{etoolbox}
\makeatletter
\providecommand{\subtitle}[1]{% add subtitle to \maketitle
  \apptocmd{\@title}{\par {\large #1 \par}}{}{}
}
\makeatother
\subtitle{― Logarithmic Transform \(O = e^{\ln A + n\theta}\),
Observer-Centered Past Light Cone, Resolution Gauge, and
\((x, y, z, t, R, Q_1, Q_2, Q_3)\) Mapping ―}
\author{Noriaki Kihara (WF System Co., Ltd.) ORCID 0009-0004-6753-4020}
\date{November 28, 2026 Version DOI 10.5281/zenodo.24625963}

\begin{document}
\maketitle

{
\setcounter{tocdepth}{2}
\tableofcontents
}
\textbf{Author:} Noriaki Kihara (WF System Co., Ltd.)\\
\textbf{ORCID:} 0009-0004-6753-4020\\
\textbf{Date:} November 28, 2026\\
\textbf{Version DOI:} 10.5281/zenodo.24625963\\
\textbf{Concept DOI:} 10.5281/zenodo.24625962\\
\textbf{Status:} Readout connection hypothesis paper (\textbf{Part III})
following the Unified Full Paper Part I (Concept DOI:
\texttt{10.5281/zenodo.24622200}) and Unified Connection Hypothesis Part
II (Concept DOI: \texttt{10.5281/zenodo.24622565}). Motivated directly
by two concrete gaps remaining after Part I and Part II---(i) the lack
of an explicit readout map from the \(\mathrm{Tri}(T)\) state space to
\((x, y, z, t, R, Q_1, Q_2, Q_3)\) and (ii) the absence of a
resolution-based definition of visibility---this paper directly connects
the readout geometry developed in prior thought experiments
{[}K9--K14{]} to the actual deterministic \(\mathrm{Tri}(T)\) states.

\begin{center}\rule{0.5\linewidth}{0.5pt}\end{center}

\subsection{Abstract}\label{abstract}

Part I {[}P1{]} established without external axioms or random numbers
that the single normalized real 3-tensor map
\(\mathrm{Tri}(T)_{ijk} = \sum_{p,q,r} T_{ipq} T_{qjr} T_{rpk}\)
spontaneously generates Hurwitz crystal selection, four quantum theorems
(Q1--Q4), spherical 4-class metric distortion with critical collapse,
Hurwitz mass hierarchy, 28D \(Spin(8)\) Triality spin-\(1/2\), and CAR
\texttt{1:3:3:1} fermion charges. Part II {[}P2{]} introduced the
\(O_8^{\mathrm{ext}} \otimes O_8^{\mathrm{int}}\) separation hypothesis
to connect gravity, electromagnetism, QCD, and Dirac matter.

Nevertheless, two essential gaps remained before closing the connection
to general relativity and the Standard Model: (1) how and why the
\(\mathrm{Tri}(T)\) state space (phase space) maps onto physical
spacetime and internal coordinates \((x, y, z, t, R, Q_1, Q_2, Q_3)\),
and (2) how ``visible vs.~hidden'' is defined by an observer's
resolution rather than by inspecting all \(64^3\) components from an
external God's-eye view. In this Part III, we bridge our prior readout
frameworks---central projection and split-octonion null balance
{[}K9{]}, two sections of the complex exponent {[}K10{]}, resolution as
a gauge {[}K11{]}, light-cone projection axes {[}K12, K13{]}, and
cross-term visibility {[}K14{]}---directly onto the deterministic
\(\mathrm{Tri}(T)\) states
(\texttt{tri\_loop\_spacetime\_log\_resolution\_bridge\_v0.py}, no
random numbers) and establish five results:

\begin{enumerate}
\def\labelenumi{\arabic{enumi}.}
\tightlist
\item
  \textbf{\texttt{(1+3)+(1+3)} Split-Octonion Null Balance \&
  Projection-Axis Covariance}: On the Part I vacuum time crystal
  \(H = O_8 \otimes M_8\), all seven Fano quaternion
  \texttt{(1+3)+(1+3)} splits strictly satisfy
  \(|a|^2 = |b|^2 = 0.500000\) (max difference
  \textbf{\texttt{3.331e-16}}), and 1D (\(\mathbb{R}\)), 2D
  (\(\mathbb{C}\)), and 4D (\(\mathbb{H}\)) projection-axis choices
  preserve the exact ratio \(d/8\) across all selections (max deviation
  \texttt{5.55e-17} to \texttt{2.22e-16}).
\item
  \textbf{\texttt{1\ +\ 3\ +\ 4} Split of Outer \(O_8^{\mathrm{ext}}\)
  and Pure \(SU(2)\)}: Inside \(G_2\) (14D), the subgroup preserving the
  \texttt{3} (\(\{e_1,e_2,e_3\}\)) and \texttt{4}
  (\(\{e_4,e_5,e_6,e_7\}\)) sectors is 6D \(SO(4)\), containing a
  \textbf{3D pure \(SU(2)\) subgroup that leaves 3-space strictly
  invariant (max error \texttt{8.327e-17}) while rotating the 4-space as
  spin-\(1/2\) (\(2\pi\) rotation \(= -I\), error \texttt{5.778e-15})}.
  Thus \(O_8^{\mathrm{ext}} \otimes O_8^{\mathrm{int}}\) (64 real DOF)
  accommodates gravity (\texttt{1+3}), electroweak \(SU(2)_L\)/Higgs
  4-component doublet (outer \texttt{4}), and \(SU(3)_C \times U(1)_Q\)
  (inner \texttt{2+6}) without overlap.
\item
  \textbf{Logarithmic Transform \(O = e^{\ln A + n\theta}\) (Origin of
  Duality \& Local Curvature Radius \(R\))}: Because \(\mathrm{Tri}\)
  operates in the multiplicative state space while spacetime and
  rapidity live in the additive exponent, \(\ln A\) and \(n\theta\) sit
  side by side, yielding Born's rule \(\cos^2(\theta/2)\) and the
  Lorentz factor \(\cosh\phi\) as two sections of one identity (error
  \texttt{2.220e-16}). Re-reading the Part I §6 mass sweep via
  \(\Delta\phi = \frac{1}{2}\ln(E_{00}/E_{\mathrm{far}})\) and
  \(R_{\mathrm{loc}}/R_0 = e^{-\Delta\phi}\) shows monotonic shrinking
  \texttt{1.0000\ \$\textbackslash{}to\$\ 0.8328} for
  \(\epsilon = 0 \to 0.270\), followed by a discontinuous horizon cut
  \(R_{\mathrm{loc}}/R_0 \to 0\) at \(\epsilon = 0.275\).
\item
  \textbf{Observer-Centered Past Light Cone \& Tidal Splitting}:
  Interpreting \(n\) not as shared Cartesian \((x,y,z)\) but as
  sightlines on the observer's own past light cone, the logarithmic past
  distance \(D_{\log}(A_{\mathrm{obs}}\to B)\) is strictly isotropic in
  vacuum (\(0.000000 \pm 10^{-16}\), equivalence principle) for any
  observer, whereas in the presence of mass P4 a distant observer
  (\(A_{\mathrm{obs}}=1\)) measures a clear \textbf{tidal splitting
  \(\Delta D_{\mathrm{tidal}} = +0.028615\)} between the radial mass
  direction (\texttt{-0.006094}) and the six tangential directions
  (\texttt{-0.034709\ ±\ 5.27e-16}).
\item
  \textbf{Resolution Gauge Definition of Visibility}: (a) Defining
  visibility at \(\tau = 75\) by \(L_3 > \delta_{\mathrm{res}}\) and
  sweeping \(\delta_{\mathrm{res}}\) across 18 orders of magnitude
  (\texttt{1e-2} to \texttt{1e-20}) leaves the visible set (64
  triangles) and weighted view angle
  \(\langle\rho\rangle_w = 0.7767719701\) invariant to 10 digits
  (resolution = gauge). (b) Under the \(\mathbb{Z}_2^3\) Walsh--Hadamard
  low-pass window \(K_{\mathrm{cut}} \in \{0,1,2,3\}\), the minimum
  resolution required to fully resolve Hurwitz particles
  \(\mathbb{O}_8, \mathbb{H}_4, \mathbb{C}_2, \mathbb{R}_1\)
  (codimension \(m=0,1,2,3\)) strictly matches \(K_{\mathrm{cut}} = m\).
\end{enumerate}

Section 6 cleanly separates the remaining open problems (multi-shell
\(O_8^{\otimes K}\) continuum limit, \(J_3(\mathbb{O})\) and three
generations, dynamical electroweak breaking and parity violation,
\(\alpha \approx 1/137\), two-particle dynamical scattering, and
dynamical CHSH violation).

\begin{center}\rule{0.5\linewidth}{0.5pt}\end{center}

\subsection{1. Introduction: Three Remaining Dissatisfactions as the
Motivation
Chain}\label{introduction-three-remaining-dissatisfactions-as-the-motivation-chain}

In a paper centered on physical readout and interpretation, stating
conclusions first invites the impression of post-hoc fitting. The
motivation for Part III is the direct chain of three concrete
dissatisfactions remaining after Part I {[}P1{]} and Part II {[}P2{]}:

\begin{enumerate}
\def\labelenumi{\arabic{enumi}.}
\tightlist
\item
  \textbf{Dissatisfaction 1 (Unclosed Spacetime Connection)}: Although
  Part I §6 derived 4-class metric distortion and critical collapse and
  Part II connected gravity, EM, QCD, and Dirac matter, our prior work
  established central projection
  \(\sum z_n^2 = 0 \mapsto \sum z_n^2 = R^2\) {[}K9, K12{]} and the 8D
  sign-vector null equation
  \[x^2 + y^2 + z^2 + (it)^2 = R^2 + Q_1^2 + Q_2^2 + Q_3^2. \tag{1}\]
  Part I and Part II lacked the explicit bridge explaining how and why
  the \(\mathrm{Tri}(T)\) state space maps onto
  \((x, y, z, t, R, Q_1, Q_2, Q_3)\).
\item
  \textbf{Dissatisfaction 2 (Missing Resolution Definition of
  Visibility)}: Part I §2.3 measured the view angle \(\rho = L_3/L_1\)
  by inspecting all \(64^3\) tensor entries externally. What was missing
  is the principle from {[}K11, K14{]}: observation is not ``existence
  vs.~non-existence'' but ``visible vs.~hidden relative to an observer's
  resolution,'' where resolution is a gauge freedom and the view angle
  is its invariant.
\item
  \textbf{Dissatisfaction 3 (Misreading \(n\) as Shared \(xyz\) and the
  Role of \texttt{1–7} Octonion Sectors)}: Reading the unit imaginary
  vector \(n\) in \(O = e^{\ln A + n\theta}\) as shared Cartesian
  \((x,y,z)\) violates background independence; the apex of the light
  cone is always the observer, so \(n\) labels sightlines on the
  observer's \textbf{past light cone (sector coordinates)}. Moreover,
  when reading \texttt{1–7} as \texttt{(1+3)+(1+3)} or \texttt{1+3+4}
  inside \(O_8^{\mathrm{ext}} \otimes O_8^{\mathrm{int}}\), the physical
  role of the extra \texttt{4} components must be identified.
\end{enumerate}

As in Part II, every claim is explicitly tagged as \textbf{【Known】},
\textbf{【Part I/II Established】}, \textbf{【Part III Verified】}, or
\textbf{【Hypothesis / Unverified】}.

\begin{center}\rule{0.5\linewidth}{0.5pt}\end{center}

\subsection{\texorpdfstring{2. \texttt{(1+3)+(1+3)} and \texttt{1+3+4}
Readings of \(O_8\) and the 64D Role
Table}{2. (1+3)+(1+3) and 1+3+4 Readings of O\_8 and the 64D Role Table}}\label{and-134-readings-of-o_8-and-the-64d-role-table}

\subsubsection{2.1 Split-Octonion Null Balance and 1D/2D/4D
Projection-Axis Covariance 【Part III Verified
P3-1】}\label{split-octonion-null-balance-and-1d2d4d-projection-axis-covariance-part-iii-verified-p3-1}

Testing the actual Part I vacuum time crystal \(H = O_8 \otimes M_8\)
across all seven Fano quaternion subalgebras \(H_4 \subset O_8\) and
across 1D (\(\mathbb{R}\)), 2D (\(\mathbb{C}\)), and 4D (\(\mathbb{H}\))
projection axes {[}K9, K13{]} yields:

\begin{itemize}
\tightlist
\item
  \textbf{All 7 \texttt{(1+3)+(1+3)} splits}:
  \(|a|^2 = |b|^2 = 0.500000\) with max difference
  \textbf{\texttt{3.331e-16}} (exact split-octonion null balance).
\item
  \textbf{1D/2D/4D projection-axis covariance}: 1D (\(\mathbb{R}\),
  \(1+7\)) = \texttt{0.125000} (dev \texttt{5.551e-17}), 2D
  (\(\mathbb{C}\), \(2+6\)) = \texttt{0.250000} (dev
  \texttt{2.220e-16}), 4D (\(\mathbb{H}\), \(4+4\)) = \texttt{0.500000}
  (dev \texttt{2.220e-16}).
\item
  \textbf{With localized mass P4 (\(\epsilon=0.16, \tau=40\))}:
  \(|a|^2 = 0.499002\) vs.~\(|b|^2 = 0.499488\) (ratio
  \texttt{0.999026}, \(\Delta\phi_{ab} = -0.000487\)), reflecting the
  greater multiplicity of tangential edges (\texttt{+7.2\%}) relative to
  the central site (\texttt{+32.5\%}).
\end{itemize}

\subsubsection{\texorpdfstring{2.2 \texttt{1\ +\ 3\ +\ 4} Split of Outer
\(O_8^{\mathrm{ext}}\): Pure \(SU(2)\) Subgroup 【Part III Verified
P3-5】}{2.2 1 + 3 + 4 Split of Outer O\_8\^{}\{\textbackslash mathrm\{ext\}\}: Pure SU(2) Subgroup 【Part III Verified P3-5】}}\label{split-of-outer-o_8mathrmext-pure-su2-subgroup-part-iii-verified-p3-5}

A single \(O_8\) (8 real DOF) lacks sufficient degrees of freedom for
all Standard Model quantum numbers; Dixon and Furey's
\(\mathbb{C}\otimes\mathbb{H}\otimes\mathbb{O}\) likewise has
\(2\times 4\times 8 = 64\) real DOF {[}E1, E2{]}. Extracting via SVD the
subgroup of \(G_2\) (14D) on \(O_8^{\mathrm{ext}}\) that preserves the
\texttt{3} (\(\{e_1,e_2,e_3\}\)) and \texttt{4}
(\(\{e_4,e_5,e_6,e_7\}\)) sectors yields:

\begin{itemize}
\tightlist
\item
  \textbf{Block-preserving subgroup}: strictly \textbf{6D
  \(SO(4) \cong SU(2)_{\mathrm{space}} \times SU(2)_{\mathrm{pure4}}\)}.
\item
  \textbf{Pure 4-space subgroup}: \textbf{3D
  \(SU(2)_{\mathrm{pure4}}\)}, whose action on 3-space
  \(\{e_1,e_2,e_3\}\) vanishes to \textbf{\texttt{8.327e-17}}.
\item
  \textbf{Spin-\(1/2\) doublet transformation}: under a \(2\pi\)
  rotation by \(SU(2)_{\mathrm{pure4}}\), the 4-space
  \(\{e_4,e_5,e_6,e_7\}\) flips sign to \(-I\) with error
  \textbf{\texttt{5.778e-15}}, matching the geometric carrier of
  \textbf{weak isospin \(SU(2)_L\) / the 4-real-component Higgs doublet}
  【Hypothesis / Unverified】.
\end{itemize}

Table 1 (and Fig. 2(b)) summarizes the complete, non-overlapping 64D
assignment of \(O_8^{\mathrm{ext}} \otimes O_8^{\mathrm{int}}\).

{\def\LTcaptype{none} % do not increment counter
\begin{longtable}[]{@{}
  >{\raggedright\arraybackslash}p{(\linewidth - 8\tabcolsep) * \real{0.2000}}
  >{\raggedright\arraybackslash}p{(\linewidth - 8\tabcolsep) * \real{0.2000}}
  >{\raggedright\arraybackslash}p{(\linewidth - 8\tabcolsep) * \real{0.2000}}
  >{\raggedright\arraybackslash}p{(\linewidth - 8\tabcolsep) * \real{0.2000}}
  >{\raggedright\arraybackslash}p{(\linewidth - 8\tabcolsep) * \real{0.2000}}@{}}
\toprule\noalign{}
\begin{minipage}[b]{\linewidth}\raggedright
Sector
\end{minipage} & \begin{minipage}[b]{\linewidth}\raggedright
Split
\end{minipage} & \begin{minipage}[b]{\linewidth}\raggedright
Exact Symmetry
\end{minipage} & \begin{minipage}[b]{\linewidth}\raggedright
Physical DOF / Coordinates
\end{minipage} & \begin{minipage}[b]{\linewidth}\raggedright
Status
\end{minipage} \\
\midrule\noalign{}
\endhead
\bottomrule\noalign{}
\endlastfoot
\textbf{Outer \(O_8^{\mathrm{ext}}\)} & \textbf{\texttt{1}} (\(e_0\)) &
Scalar (\(\ln A\)) & Time reference \(t\), scale, local curvature radius
\(R\) & 【Part III Verified】+【Hypothesis】 \\
& \textbf{\texttt{3}} (\(e_1,e_2,e_3\)) & \(SO(3)\) (3D) & Past
light-cone sightlines \(n\) (\(x,y,z\) readout), gravity tides & 【Part
I/III Verified】 \\
& \textbf{\texttt{4}} (\(e_4,e_5,e_6,e_7\)) & \textbf{Pure \(SU(2)\)
(3D)} & \textbf{Weak isospin \(SU(2)_L\), Higgs 4-comp. doublet,
spin-\(1/2\)} & 【Part III Verified】+【Hypothesis】 \\
\textbf{Inner \(O_8^{\mathrm{int}}\)} & \textbf{\texttt{2\ +\ 6}}
(\(e_7\) fixed) & \textbf{\(SU(3)_C\) (8D) \(\times U(1)_Q\)} &
\textbf{Strong/EM forces, CAR \texttt{1:3:3:1} charges
\((Q_1,Q_2,Q_3)\), clock \(M_8\)} & 【Part I/II Established】 \\
\end{longtable}
}

\begin{center}\rule{0.5\linewidth}{0.5pt}\end{center}

\subsection{\texorpdfstring{3. Logarithmic Transform
\(O = e^{\ln A + n\theta}\): Origin of Duality and Local Curvature
Radius
\(R\)}{3. Logarithmic Transform O = e\^{}\{\textbackslash ln A + n\textbackslash theta\}: Origin of Duality and Local Curvature Radius R}}\label{logarithmic-transform-o-eln-a-ntheta-origin-of-duality-and-local-curvature-radius-r}

\subsubsection{3.1 Multiplicative Dynamics vs.~Additive Spacetime
【Known】+【Part III Verified
P3-2】}\label{multiplicative-dynamics-vs.-additive-spacetime-knownpart-iii-verified-p3-2}

Any octonion \(O \in \mathbb{O}\) admits the single-exponent form
\[O = A\,e^{n\theta} = e^{\ln A + n\theta} \tag{2}\] (round-trip error
\texttt{1.665e-16}). While \(\mathrm{Tri}(T)\) evolves in the
\textbf{multiplicative} state space, physical spacetime intervals,
rapidities, and energies live in the \textbf{additive} exponent
(\(\ln A + n\theta\)). Because the logarithm places scale
\(\phi = \ln A\) and phase \(n\theta\) in the same linear space, the
identity \(|\cosh(w/2)|^2 \pm |\sinh(w/2)|^2 = (\cosh\phi, \cos\theta)\)
(error \textbf{\texttt{2.220e-16}}) {[}K10{]} renders Born's rule
(\(\cos^2(\theta/2)\) at \(\phi=0\)) and the Lorentz factor
(\(\cosh\phi\) at \(\theta=0\)) dual sections related by a \(90^\circ\)
rotation. Normalization by Frobenius norm corresponds merely to a
uniform origin shift in \(\ln A\), and a compact bounded state space
{[}K14{]} maps to an open horizon
(\(A \to 0 \Rightarrow \ln A \to -\infty\)).

\subsubsection{\texorpdfstring{3.2 Local Curvature Radius
\(R_{\mathrm{loc}}\) Shrinking and Horizon Cut 【Part III Verified
P3-2】}{3.2 Local Curvature Radius R\_\{\textbackslash mathrm\{loc\}\} Shrinking and Horizon Cut 【Part III Verified P3-2】}}\label{local-curvature-radius-r_mathrmloc-shrinking-and-horizon-cut-part-iii-verified-p3-2}

Re-reading the Part I §6 mass P4 sweep (\(\tau=40\)) via
\(\Delta\phi = \frac{1}{2}\ln(E_{00}/E_{\mathrm{far}})\) and
\(R_{\mathrm{loc}}/R_0 = e^{-\Delta\phi}\) {[}K12{]} yields Table 2 and
Fig. 1(a).

{\def\LTcaptype{none} % do not increment counter
\begin{longtable}[]{@{}
  >{\centering\arraybackslash}p{(\linewidth - 10\tabcolsep) * \real{0.1667}}
  >{\centering\arraybackslash}p{(\linewidth - 10\tabcolsep) * \real{0.1667}}
  >{\centering\arraybackslash}p{(\linewidth - 10\tabcolsep) * \real{0.1667}}
  >{\centering\arraybackslash}p{(\linewidth - 10\tabcolsep) * \real{0.1667}}
  >{\centering\arraybackslash}p{(\linewidth - 10\tabcolsep) * \real{0.1667}}
  >{\centering\arraybackslash}p{(\linewidth - 10\tabcolsep) * \real{0.1667}}@{}}
\toprule\noalign{}
\begin{minipage}[b]{\linewidth}\centering
Mass amplitude \(\epsilon\)
\end{minipage} & \begin{minipage}[b]{\linewidth}\centering
\(E_{00}/E_{\mathrm{far}}\)
\end{minipage} & \begin{minipage}[b]{\linewidth}\centering
\(\Delta\phi = \frac{1}{2}\ln(E_{00}/E_{\mathrm{far}})\)
\end{minipage} & \begin{minipage}[b]{\linewidth}\centering
\textbf{\(R_{\mathrm{loc}}/R_0 = e^{-\Delta\phi}\)}
\end{minipage} & \begin{minipage}[b]{\linewidth}\centering
\(\cosh(\Delta\phi)\)
\end{minipage} & \begin{minipage}[b]{\linewidth}\centering
\(E_{\mathrm{far}}\)
\end{minipage} \\
\midrule\noalign{}
\endhead
\bottomrule\noalign{}
\endlastfoot
\texttt{0.000} & \texttt{1.000000} & \texttt{0.000000} &
\textbf{\texttt{1.000000}} & \texttt{1.000000} & \texttt{0.015625} \\
\texttt{0.050} & \texttt{1.047528} & \texttt{+0.023217} &
\textbf{\texttt{0.977051}} & \texttt{1.000270} & \texttt{0.015495} \\
\texttt{0.100} & \texttt{1.174017} & \texttt{+0.080216} &
\textbf{\texttt{0.922917}} & \texttt{1.003219} & \texttt{0.015162} \\
\texttt{0.160} & \texttt{1.324838} & \texttt{+0.140645} &
\textbf{\texttt{0.868798}} & \texttt{1.009907} & \texttt{0.014789} \\
\texttt{0.220} & \texttt{1.400375} & \texttt{+0.168370} &
\textbf{\texttt{0.845041}} & \texttt{1.014208} & \texttt{0.014686} \\
\texttt{0.270} & \texttt{1.441867} & \texttt{+0.182969} &
\textbf{\texttt{0.832794}} & \texttt{1.016786} & \texttt{0.014492} \\
\textbf{\texttt{0.275}} & \textbf{\(\infty\) (horizon cut)} &
\textbf{\(+\infty\)} & \textbf{\texttt{0.000000}} & \textbf{\(\infty\)}
& \textbf{\texttt{0.000e+00}} \\
\end{longtable}
}

This directly connects the working hypothesis of {[}K12{]} (``gravity as
shrinking of the local curvature radius \(R\) and event horizon'') with
the Part I critical collapse of \(\mathrm{Tri}(T)\).

\begin{center}\rule{0.5\linewidth}{0.5pt}\end{center}

\subsection{4. Observer-Centered Past Light Cone and Tidal Splitting
【Part III Verified
P3-3】}\label{observer-centered-past-light-cone-and-tidal-splitting-part-iii-verified-p3-3}

Treating \(n\) as sightlines from an observer site \(A_{\mathrm{obs}}\)
on its own past light cone, we define the logarithmic past distance to
site \(B \neq A_{\mathrm{obs}}\) by
\[D_{\log}(A_{\mathrm{obs}} \to B) \equiv -\frac{1}{2}\ln\!\left(\frac{E_{A_{\mathrm{obs}}, B}}{E_{A_{\mathrm{obs}}, A_{\mathrm{obs}}}}\right). \tag{3}\]

\begin{itemize}
\tightlist
\item
  \textbf{Vacuum \(H\) (equivalence principle)}: For both
  \(A_{\mathrm{obs}}=0\) and \(A_{\mathrm{obs}}=1\), all seven
  directions are strictly isotropic:
  \(D_{\log} = 0.000000 \pm 8.32\times 10^{-16}\) (and
  \(\pm 2.29\times 10^{-11}\) at \(A=1\)). Uniform curvature is
  invisible; the local readout is flat {[}K9{]}.
\item
  \textbf{With mass P4 (\(A=0, \epsilon=0.16\))}:

  \begin{itemize}
  \tightlist
  \item
    Observer on the mass (\(A_{\mathrm{obs}}=0\)): all 7 directions
    remain strictly spherical
    (\(D_{\log} = +0.134551 \pm 7.77\times 10^{-16}\)).
  \item
    \textbf{Distant observer (\(A_{\mathrm{obs}}=1\))}: the radial mass
    direction (\(B=0\)) gives \texttt{D\_log\ =\ -0.006094}, whereas the
    six tangential directions (\(B=2\dots 7\)) give
    \texttt{D\_log\ =\ -0.034709\ ±\ 5.27e-16}, producing a clear
    \textbf{tidal splitting \(\Delta D_{\mathrm{tidal}} = +0.028615\)}
    (Fig. 1(b)).
  \end{itemize}
\end{itemize}

This directly verifies the prediction of {[}K9{]} that uniform gravity
disappears in an observer-centered readout while directional tidal
distortion survives.

\begin{figure}[H]
\centering
\includegraphics[width=0.96\linewidth,height=\textheight,keepaspectratio,alt={Part III Verification (1): (a) Local curvature radius shrinking R\_\{\textbackslash mathrm\{loc\}\}/R\_0 = e\^{}\{-\textbackslash Delta\textbackslash phi\} and critical horizon cut via the logarithmic map; (b) Observer-centered past light cone showing vacuum isotropy (equivalence principle) vs.~distant-observer tidal splitting \textbackslash Delta D\_\{\textbackslash mathrm\{tidal\}\} = +0.0286.}]{figs_en/fig_p3_log_curvature_and_tidal_en.png}
\caption{Part III Verification (1): (a) Local curvature radius shrinking
\(R_{\mathrm{loc}}/R_0 = e^{-\Delta\phi}\) and critical horizon cut via
the logarithmic map; (b) Observer-centered past light cone showing
vacuum isotropy (equivalence principle) vs.~distant-observer tidal
splitting \(\Delta D_{\mathrm{tidal}} = +0.0286\).}
\end{figure}

\begin{center}\rule{0.5\linewidth}{0.5pt}\end{center}

\subsection{5. Resolution-Gauge Definition of Visibility 【Part III
Verified
P3-4】}\label{resolution-gauge-definition-of-visibility-part-iii-verified-p3-4}

To replace the external God's-eye view in Part I, we implement two
resolution filters {[}K11, K14{]} directly on \(\mathrm{Tri}(T)\) states
(Fig. 2(a)):

\subsubsection{5.1 18-Order Threshold Sweep: Resolution as Gauge, View
Angle as
Invariant}\label{order-threshold-sweep-resolution-as-gauge-view-angle-as-invariant}

At \(\tau = 75\) (representative particle component \(i=4, j=0, k=4\)),
defining visibility by \textbf{\(L_3 > \delta_{\mathrm{res}}\)} and
sweeping
\(\delta_{\mathrm{res}} \in \{10^{-2}, 10^{-5}, 10^{-10}, 10^{-20}\}\)
leaves the visible set strictly at \textbf{64 triangles} and the
weighted view angle at \textbf{\(\langle\rho\rangle_w = 0.7767719701\)}
(10-digit invariant). Detached background triangles fall below
resolution (\texttt{null} {[}K14{]}), confirming that resolution is a
gauge and the view angle is a gauge invariant {[}K11{]}.

\subsubsection{\texorpdfstring{5.2 \(\mathbb{Z}_2^3\) Harmonic Low-Pass
Window \(K_{\mathrm{cut}}\) and Hurwitz Visibility
Boundary}{5.2 \textbackslash mathbb\{Z\}\_2\^{}3 Harmonic Low-Pass Window K\_\{\textbackslash mathrm\{cut\}\} and Hurwitz Visibility Boundary}}\label{mathbbz_23-harmonic-low-pass-window-k_mathrmcut-and-hurwitz-visibility-boundary}

Applying the \(\mathbb{Z}_2^3\) Walsh--Hadamard low-pass window
\(K_{\mathrm{cut}} \in \{0,1,2,3\}\) to the exact Part I §7 Hurwitz
4-level probes (\(\epsilon=0.05, \tau=30\)) yields Table 3.

{\def\LTcaptype{none} % do not increment counter
\begin{longtable}[]{@{}
  >{\raggedright\arraybackslash}p{(\linewidth - 10\tabcolsep) * \real{0.1379}}
  >{\centering\arraybackslash}p{(\linewidth - 10\tabcolsep) * \real{0.1724}}
  >{\centering\arraybackslash}p{(\linewidth - 10\tabcolsep) * \real{0.1724}}
  >{\centering\arraybackslash}p{(\linewidth - 10\tabcolsep) * \real{0.1724}}
  >{\centering\arraybackslash}p{(\linewidth - 10\tabcolsep) * \real{0.1724}}
  >{\centering\arraybackslash}p{(\linewidth - 10\tabcolsep) * \real{0.1724}}@{}}
\toprule\noalign{}
\begin{minipage}[b]{\linewidth}\raggedright
Hurwitz Particle (codimension \(m\))
\end{minipage} & \begin{minipage}[b]{\linewidth}\centering
\(K_{\mathrm{cut}}=0\) (DC)
\end{minipage} & \begin{minipage}[b]{\linewidth}\centering
\(K_{\mathrm{cut}}\le 1\)
\end{minipage} & \begin{minipage}[b]{\linewidth}\centering
\(K_{\mathrm{cut}}\le 2\)
\end{minipage} & \begin{minipage}[b]{\linewidth}\centering
\(K_{\mathrm{cut}}\le 3\) (Full)
\end{minipage} & \begin{minipage}[b]{\linewidth}\centering
Full-visibility \(k_{\max}\)
\end{minipage} \\
\midrule\noalign{}
\endhead
\bottomrule\noalign{}
\endlastfoot
\textbf{\(\mathbb{O}_8\) (\(m=0\): background)} & \texttt{1.000000} &
\texttt{1.000000} & \texttt{1.000000} & \texttt{1.000000} &
\textbf{\texttt{0}} \\
\textbf{\(\mathbb{H}_4\) (\(m=1\): P5)} & \texttt{0.900000} &
\texttt{1.000000} & \texttt{1.000000} & \texttt{1.000000} &
\textbf{\texttt{1}} \\
\textbf{\(\mathbb{C}_2\) (\(m=2\))} & \texttt{0.749250} &
\texttt{0.916417} & \texttt{1.000000} & \texttt{1.000000} &
\textbf{\texttt{2}} \\
\textbf{\(\mathbb{R}_1\) (\(m=3\): P4)} & \texttt{0.556705} &
\texttt{0.746689} & \texttt{0.936672} & \texttt{1.000000} &
\textbf{\texttt{3}} \\
\end{longtable}
}

The codimension \(m\) strictly equals the minimum resolution window
\(K_{\mathrm{cut}} = m\) needed to fully resolve each particle; for
\(K_{\mathrm{cut}} < m\), finer modes (\(k > K_{\mathrm{cut}}\)) remain
hidden as internal rigid-crystal structure {[}K11{]}.

\begin{figure}[H]
\centering
\includegraphics[width=0.96\linewidth,height=\textheight,keepaspectratio,alt={Part III Verification (2): (a) Hurwitz 4-level particle visibility boundary at K\_\{\textbackslash mathrm\{cut\}\} = m and 18-order threshold invariance of the view angle; (b) Complete 64D readout mapping from O\_8\^{}\{\textbackslash mathrm\{ext\}\} \textbackslash otimes O\_8\^{}\{\textbackslash mathrm\{int\}\} to (x, y, z, t, R, Q\_1, Q\_2, Q\_3) and SU(2)\_L.}]{figs_en/fig_p3_resolution_gauge_and_mapping_en.png}
\caption{Part III Verification (2): (a) Hurwitz 4-level particle
visibility boundary at \(K_{\mathrm{cut}} = m\) and 18-order threshold
invariance of the view angle; (b) Complete 64D readout mapping from
\(O_8^{\mathrm{ext}} \otimes O_8^{\mathrm{int}}\) to
\((x, y, z, t, R, Q_1, Q_2, Q_3)\) and \(SU(2)_L\).}
\end{figure}

\begin{center}\rule{0.5\linewidth}{0.5pt}\end{center}

\subsection{6. Clean Separation of Remaining Open
Problems}\label{clean-separation-of-remaining-open-problems}

The following essential problems remain \textbf{【Hypothesis /
Unverified】} and are explicitly deferred to subsequent work:

\begin{enumerate}
\def\labelenumi{\arabic{enumi}.}
\tightlist
\item
  \textbf{Multi-shell \(O_8^{\otimes K}\) continuum limit}: Deriving
  \(1/r\) Newtonian gravity, \(1/r^2\) Coulomb force, light-cone
  propagation speed \(c\), Einstein's field equations, and cumulative
  gravitational time dilation \(g_{00}(r)\) by interpreting nested
  shells \(O_8^{\otimes K}\) as concentric depths of the past light
  cone.
\item
  \textbf{Dynamical electroweak breaking and maximal parity violation}:
  Dynamically coupling the outer pure \(SU(2)\) (4-space) with the inner
  \(M_8\) chiral asymmetry (\texttt{-1.41\%}), left-handed weak
  selection, and \(W/Z\) mass generation.
\item
  \textbf{Three generations, precision mass ratios, and coupling
  constants}: Connecting the exceptional Jordan algebra
  \(J_3(\mathbb{O})\) {[}E3{]} or \(O_8^{\otimes K}\) nesting to three
  fermion generations, exact lepton/quark mass ratios, and
  first-principles derivation of \(\alpha \approx 1/137.036\) and
  \(\alpha_s\).
\item
  \textbf{Dynamical scattering and multi-body quantum correlations}:
  Two-particle attractive/repulsive scattering, electron--photon
  scattering, and dynamical CHSH Bell inequality violation.
\end{enumerate}

\begin{center}\rule{0.5\linewidth}{0.5pt}\end{center}

\subsection{7. Reproducibility and Published
Files}\label{reproducibility-and-published-files}

\begin{itemize}
\tightlist
\item
  Verification script:
  \texttt{tri\_loop\_spacetime\_log\_resolution\_bridge\_v0.py}
  (deterministic, \textasciitilde23 s)
\item
  Numerical log:
  \texttt{results/tri\_loop\_spacetime\_log\_resolution\_bridge\_v0.txt}
\item
  Figure script: \texttt{plot\_part3\_figures\_v0.py}
\item
  Figures:
  \texttt{results/fig\_p3\_log\_curvature\_and\_tidal\_\{ja,en\}.png},
  \texttt{results/fig\_p3\_resolution\_gauge\_and\_mapping\_\{ja,en\}.png}
\end{itemize}

\begin{center}\rule{0.5\linewidth}{0.5pt}\end{center}

\subsection{References}\label{references}

\subsubsection{A. Self-Citations}\label{a.-self-citations}

\begin{itemize}
\tightlist
\item
  \textbf{{[}P1{]}} N. Kihara, \emph{Spontaneous Emergence of Quantum,
  Gauge, Gravitational, and Fermionic Structures from the 3D
  Triangle-Loop Product (Unified Full Paper, Part I)}, Zenodo
  (2026-11-28), Concept DOI: \texttt{10.5281/zenodo.24622200}.
\item
  \textbf{{[}P2{]}} N. Kihara, \emph{Unified Connection Hypothesis for
  Gravity, Electromagnetism, QCD, and Dirac Matter in the 3D
  Triangle-Loop Product (Part II)}, Zenodo (2026-11-28), Concept DOI:
  \texttt{10.5281/zenodo.24622565}.
\item
  \textbf{{[}K9{]}} N. Kihara, \emph{Relational Readout of Quaternions
  and Inflation as Crystallization}, Zenodo (2026-10-21), Concept DOI:
  \texttt{10.5281/zenodo.23523237}.
\item
  \textbf{{[}K10{]}} N. Kihara, \emph{Did We Properly Understand Complex
  Numbers?}, Zenodo (2026-10-23), Concept DOI:
  \texttt{10.5281/zenodo.23585436}.
\item
  \textbf{{[}K11{]}} N. Kihara, \emph{Observer Resolution Is a Gauge},
  Zenodo (2026-11-10), Concept DOI: \texttt{10.5281/zenodo.24097611}.
\item
  \textbf{{[}K12{]}} N. Kihara, \emph{Emergence of Spacetime and
  Elementary Particles from a Bounded Hilbert Space and Octonions},
  Zenodo (2026-11-16), Concept DOI: \texttt{10.5281/zenodo.24267347}.
\item
  \textbf{{[}K13{]}} N. Kihara, \emph{Projection-Axis Selection and
  Covariance of Conservation Laws}, Zenodo (2026-11-16), Concept DOI:
  \texttt{10.5281/zenodo.24273723}.
\item
  \textbf{{[}K14{]}} N. Kihara, \emph{The Coin-and-Cup Thought
  Experiment}, Zenodo (2026-11-26), Concept DOI:
  \texttt{10.5281/zenodo.24562867}.
\end{itemize}

\subsubsection{B. External Primary
Literature}\label{b.-external-primary-literature}

\begin{itemize}
\tightlist
\item
  \textbf{{[}E1{]}} G. M. Dixon, \emph{Division Algebras: Octonions,
  Quaternions, Complex Numbers and the Algebraic Design of Physics},
  Kluwer Academic Publishers (1994).
\item
  \textbf{{[}E2{]}} C. Furey,
  \emph{\(SU(3)_C \times SU(2)_L \times U(1)_Y (\times U(1)_X)\) as a
  symmetry of division algebraic ladder operators}, Eur. Phys. J. C
  \textbf{78}, 375 (2018), arXiv:1806.00612.
\item
  \textbf{{[}E3{]}} J. C. Baez, \emph{The Octonions}, Bull. Amer. Math.
  Soc. \textbf{39}(2), 145--205 (2002), arXiv:math/0105155.
\item
  \textbf{{[}E4{]}} S. Coleman and J. Mandula, \emph{All Possible
  Symmetries of the \(S\) Matrix}, Phys. Rev.~\textbf{159}(5),
  1251--1256 (1967).
\end{itemize}

\end{document}
